% =============================================================================
%  Supplementary Section S5: Change within one patient.
%  Order: setting; why the shared constants cancel in a difference; precision
%  and draws per window (stab:s5windows); simulated pairs (stab:s5pairs); why
%  the ratio fails; conditions (stab:s5conditions); other instruments and the
%  saturable-clearance model; what would test the result.
% =============================================================================
\section{Change within one patient}
\label{si:s5}

Supplementary Section~\ref{si:s4} showed that no available measurement fixes the level of
residual capacity $\theta$ at 5\% of normal. This section asks a different question: whether
capacity has changed between two periods of one patient's record. We explain why the constants
that block the level cancel in a difference (Section~\ref{si:s5:cancel}), give the precision and
its dependence on the number of draws (Section~\ref{si:s5:precision}), report fits to simulated
pairs of records (Section~\ref{si:s5:pairs}), explain why the change must be read as an amount
rather than a percentage (Section~\ref{si:s5:ratio}), and list the conditions under which the
result holds (Sections~\ref{si:s5:conditions}--\ref{si:s5:test}).

% -----------------------------------------------------------------------------
\subsection{Setting}
\label{si:s5:setting}

The setting is that of Supplementary Section~\ref{si:s4:setting}: the linear-clearance model, in
which each settled draw satisfies
$(\theta+\Delta_i)\,g(m_i)+k_e m_i=L_i+b$, with $m_i$ the marker level, $\Delta_i$ the delivered
increment, $g(m)=m/(K_m+m)$, $k_e$ the rate constant of the non-enzymatic route, $L_i$ the load
recorded for that day and $b$ the unrecorded load offset. The patient starts at 5\% of normal,
$\theta_1=150$\um, on the assumed fortnightly dose train ($D=6000$\um, half-life 3.5\,d), with a
10\% day-to-day variation in intake, a 5\% assay coefficient of variation and a 5\% load-record
error, each draw at its own noise level. Each estimate comes from a window of 104 draws at
evenly spaced times over two years, one every 7.09 days, so that successive draws move through
the dose cycle. Within a window $\theta$ is taken as constant; between the two windows it may
change, from $\theta_1$ to $\theta_2$. The half-saturation $K_m$, the offset $b$ and the rate
constant $k_e$ are patient constants, shared by the two windows, while each window has its own
intake record, its own draws and its own noise. The quantity of interest is the difference
$\hat\theta_2-\hat\theta_1$, in \um.
% src: audit/brief_2026-10-01/t56_contrast_mc.py:1-20 (two windows of 104 draws over 730 d, D = 6000, 10% load variation, per-draw noise; K_m, b, k_e shared), :31,61 (N 104, WIN 730, linspace: 730/103 = 7.09 d); manuscript/main.tex:778-787

% -----------------------------------------------------------------------------
\subsection{Why the shared constants cancel in a difference}
\label{si:s5:cancel}

Read at the observed level $m_i$, an error $\delta k_e$ in the rate constant or $\delta b$ in
the load offset is absorbed by a shift in the capacity estimate of
\begin{equation}
\label{seq:s5shift}
\delta\hat\theta\;=\;-\,\delta k_e\,(K_m+m_i)\;+\;\delta b\,\frac{K_m+m_i}{m_i}.
\end{equation}
The shift is an amount, in \um, set by the marker level and not by $\theta$. On therapy the
delivered enzyme dominates the enzyme route (at the assumed dose $\theta+\Delta$ is about 14
times $\theta$ at 5\% of normal), so the two windows visit nearly the same marker levels and
their estimates are shifted by nearly the same amount. The shift then cancels in the difference
$\hat\theta_2-\hat\theta_1$ whatever the change. The half-saturation $K_m$, also shared, enters
both windows in nearly the same way for the same reason.
% src: REPORT.md:1206 (d theta = -d k_e (K_m + m) at a fixed level; additive shifts cancel in the difference); out/t56_contrast_mc.txt:18-21; out/t25_drug_cv_marker.txt:5-9 (theta + Delta = 14.0 x theta at the booked dose)

In the logarithm of the ratio the shift does not cancel. A common shift $\delta$ moves the log
ratio by about
\begin{equation}
\label{seq:s5logshift}
\ln\hat\theta_2-\ln\hat\theta_1\;\approx\;\delta\Bigl(\frac{1}{\theta_2}-\frac{1}{\theta_1}\Bigr),
\end{equation}
which vanishes only when nothing has changed. At 5\% of normal, moreover, $\delta$ is of the
order of $\theta$ itself, because the record does not bound the level from below, so the
first-order statement~\eqref{seq:s5logshift} does not hold and the log ratio has heavy tails
(Section~\ref{si:s5:ratio}).
% src: REPORT.md:1206 (common shift moves ln ratio by about Delta (1/theta2 - 1/theta1)); out/t56_contrast_mc.txt:14-21

The multiplicative constants of the direct readings behave the other way. The scale $\kappa$ of
an ex vivo assay and the recovery $c$ of a breath test multiply the reading, so they cancel in
the ratio of two readings rather than in their difference. Each patient constant cancels on the
scale on which it acts, provided it is the same in both windows. A change within one patient
therefore escapes the obstacles of Supplementary Section~\ref{si:s4}, because the constant that
each route to the level would have to import is a patient constant.
% src: REPORT.md:573 (one mechanism: additive in the difference, multiplicative in the ratio), :30

% -----------------------------------------------------------------------------
\subsection{Precision and the number of draws}
\label{si:s5:precision}

Linearised at 5\% of normal, the Fisher information of two windows of 104 draws gives a
standard deviation of 0.133 for $\ln\theta_2-\ln\theta_1$, which at $\theta=150$\um\ is 20.0\um\
for $\theta_2-\theta_1$. Relative to the patient's capacity this is 0.13, below the target of
0.20; it is the only quantity in our analysis at 5\% of normal that reaches the target from the
marker alone. Below, Fisher values of the contrast are given in this relative form; multiplied by
150 they give the standard deviation of the difference in \um.
% src: out/t56_contrast_mc.txt:4 (median 0.133; 20.0 uM/d); out/t22_longitudinal.txt:24-28; REPORT.md:571,1206

The precision is set by the number of draws in each window. With draws spread over two years, 25
draws per window give 0.26, 52 give 0.19, 104 give 0.13 and 365 give 0.07. At a fixed number of
draws the calendar span of the window barely matters (Supplementary Table~\ref{stab:s5windows}),
and about 45 draws per window reach 0.20. Two short, dense windows, such as six months of daily
draws (0.10), therefore serve at least as well as two years of weekly draws (0.13), and a longer
gap between the windows only allows more drift (Section~\ref{si:s5:conditions}). Short windows also limit the
bias from assuming $\theta$ constant within a window, which grows with the span of the window
rather than with its number of draws.
% src: out/t22_longitudinal.txt:17-25 (0.264, 0.191, 0.133, 0.067); out/t24_contrast_stress.txt:5-13 (span vs draws; ~45 draws; 6-month daily 0.096); REPORT.md:583,598-600 (drift bias -rT/2 scales with span)

\begin{table}[htbp]
\centering\small
\caption{\label{stab:s5windows}Fisher standard deviation of the within-patient contrast,
$\ln\theta_2-\ln\theta_1$ (equal to $\operatorname{sd}(\hat\theta_2-\hat\theta_1)/150$), at 5\%
of normal in the linear-clearance model. (a)~Against the number of draws per window, spread over
two years. (b)~Against the calendar span of each window and the sampling frequency, with $n$ the
number of draws per window; panel~(b) has 24 monthly draws in two years, hence 0.281 against the
0.264 of 25 draws in panel~(a).}
\begin{tabular}{@{}lcccc@{}}
\toprule
\multicolumn{5}{@{}l}{(a) draws per window, over two years}\\
\midrule
draws per window & 25 & 52 & 104 & 365 \\
contrast & 0.264 & 0.191 & 0.133 & 0.067 \\
\midrule
\multicolumn{5}{@{}l}{(b) span of each window and sampling frequency}\\
\midrule
span & monthly & weekly & daily & \\
6 months & 0.568 ($n=6$) & 0.278 ($n=26$) & 0.096 ($n=180$) & \\
1 year & 0.388 ($n=12$) & 0.187 ($n=52$) & 0.067 ($n=365$) & \\
2 years & 0.281 ($n=24$) & 0.133 ($n=104$) & 0.047 ($n=730$) & \\
\bottomrule
\end{tabular}
\end{table}
% src: out/t22_longitudinal.txt:21-25; out/t24_contrast_stress.txt:5-9

% -----------------------------------------------------------------------------
\subsection{Simulated pairs of records}
\label{si:s5:pairs}

Because the Fisher value is a linearisation, and at 5\% of normal the estimate from one window
can reach zero, we fitted simulated pairs of records. For each scenario, 400 pairs were generated
with the true capacity in the second window set to a multiple of that in the first (1, 1.2, 0.8,
1.45 or 0.7), the marker recomputed at the new capacity and the patient's load scale unchanged.
Each pair was fitted jointly by least squares from several starting points, with $\theta$ free in
each window, $K_m$, $b$ and $k_e$ tied between the windows, and $\theta\ge10^{-3}$\um. As in the
Fisher calculation, each window contributes its own copy of the priors, so the shared $k_e$ and
$b$ carry effective prior standard deviations of $0.8/\sqrt2$ on $\ln k_e$ and $0.05L/\sqrt2$.
Two tests at the two-sided 5\% level were applied: a difference test, rejecting no change when
$|\hat\theta_2-\hat\theta_1|>1.96\times20.0$\um, and a log-ratio test, rejecting when
$|\ln\hat\theta_2-\ln\hat\theta_1|>1.96\times0.133$.
% src: audit/brief_2026-10-01/t56_contrast_mc.py:7-17 (footing; priors added to both tied copies: 0.8/sqrt2, 0.05L/sqrt2; theta >= 1e-3; multi-start; window 2 at ratio x theta1 with the load scale unchanged), :35 (REPS 400); out/t56_contrast_mc.txt:4-5 (tests)

The difference behaves as the Fisher value predicts (Supplementary Table~\ref{stab:s5pairs}). It
is nearly unbiased with or without a change, its standard deviation is 19.4--21.3\um\ in every
scenario, and the test on it raises false alarms in 5\% of pairs with no change. With a standard
deviation of 20\um, a change of $2.80\times20=56$\um\ either way, 37\% of the capacity of a patient
at 5\% of normal, is detected with 80\% power. A rise of 45\% (67.5\um) is detected in 91\% of
pairs and a fall of 30\% (45\um) in 64\%; changes of 20\% (30\um) are detected in about a third of
pairs. Read as the standard deviation of a log ratio, the same Fisher value would suggest 80\%
power both for a rise of 45\% and for a fall of 30\%. The simulation shows the fall detected less
often, because the record measures a change of fixed size in \um, and a 30\% fall is a smaller
amount than a 45\% rise.
% src: out/t56_contrast_mc.txt:7-12 (table), :22-24 (2.80 x 20 = 56 uM/d, 37%); REPORT.md:1196-1206; out/t24_contrast_stress.txt:52-58 (+45%/-31% from the Fisher value)

\begin{table}[htbp]
\centering\small\setlength{\tabcolsep}{4pt}
\caption{\label{stab:s5pairs}Simulated pairs of records at 5\% of normal ($\theta_1=150$\um),
linear-clearance model, two windows of 104 draws each, 400 pairs per scenario. Floor: pairs in
which at least one window's estimate lies at the lower bound, $10^{-3}$\um. Outside: fitted log
ratios outside $\pm1.96\times0.133$ of the truth, which should be 5\%. Difference: mean and
standard deviation of $\hat\theta_2-\hat\theta_1$ (\um). Rejections: share of pairs in which each
test rejects no change at the two-sided 5\% level; with no change this is the false-alarm rate.}
\begin{tabular}{@{}cccccccc@{}}
\toprule
& true change & & & \multicolumn{2}{c}{difference} & \multicolumn{2}{c}{rejections} \\
\cmidrule(lr){5-6}\cmidrule(l){7-8}
$\theta_2/\theta_1$ & (\um) & floor & outside & mean & sd & difference test & log-ratio test \\
\midrule
1.00 & 0 & 23\% & 34\% & $+1.0$ & 20.2 & 5\% & 34\% \\
1.20 & $+30$ & 20\% & 30\% & $+31.2$ & 21.3 & 37\% & 42\% \\
0.80 & $-30$ & 23\% & 34\% & $-30.3$ & 19.6 & 33\% & 45\% \\
1.45 & $+67.5$ & 24\% & 41\% & $+66.9$ & 19.9 & 91\% & 66\% \\
0.70 & $-45$ & 28\% & 48\% & $-46.0$ & 19.4 & 64\% & 61\% \\
\bottomrule
\end{tabular}
\end{table}
% src: out/t56_contrast_mc.txt:7-12; REPORT.md:1198-1204

% -----------------------------------------------------------------------------
\subsection{Why the change must be read as an amount}
\label{si:s5:ratio}

The ratio $\hat\theta_2/\hat\theta_1$ fails at 5\% of normal. In 20--28\% of pairs the estimate
of at least one window falls to the lower bound, because the record does not bound $\theta$ from below
(main text, Fig.~3c), and the shared constants then pull the other estimate towards zero as well.
The ratio of two numbers near zero can take any value. With no change, the middle of the fitted
log ratios roughly follows the Fisher value (a robust standard deviation of 0.174 against 0.133),
but their central 95\% spans $-10.0$ to $+10.3$, and 34\% of them fall outside the range that
should hold 95\%. A test on the log ratio therefore raises false alarms in about a third of
pairs, and under a true change its rejections mix detections with these false alarms. A
percentage change would need the level of $\theta$, which the record does not provide; what the
record measures is the change in \um.
% src: out/t56_contrast_mc.txt:7-8 (floor 0.20-0.28; robust sd 0.174; [-10.04, +10.33]; outside 0.34), :14-22; REPORT.md:1206

% -----------------------------------------------------------------------------
\subsection{Conditions}
\label{si:s5:conditions}

The contrast of 0.13 is a floor computed under two assumptions: that each marker value scatters
by 5\% around the model's prediction for that draw, and that the patient constants are exactly
the same in both windows. Supplementary Table~\ref{stab:s5conditions} shows how the contrast
degrades when either assumption fails, and how two events between the windows act on it.
% src: REPORT.md:605; out/t30_monitoring_confounders.txt:5-7

\paragraph{Per-draw scatter of the marker.} The 5\% is the assay's coefficient of variation. In
a patient the relevant quantity is the scatter of each value around the model's prediction for
that draw, after the predictable swings over the dose cycle and with the recorded intake are
accounted for. No such value has been published for 3-hydroxypropionate, propionylcarnitine (C3)
or the C3/C2 ratio in stable PA. The one measured serial variability of a related marker, serum
C3 within a day in healthy adults, averages 26.4\% (range 8.2--61.1\%)~\cite{thompson2012c3}. At
that value the contrast is 0.475, about 71\um, and about 590 draws per window would be needed to
reach 0.20. In methylmalonic acidaemia the between-visit coefficient of variation of serum
methylmalonic acid is higher still, $31\pm25\%$~\cite{breath_test_mma}. The healthy-adult value is
a weak anchor in both directions: it was measured within one day, without standardised sampling,
on a small and labile plasma pool, which may overstate the scatter in a standardised clinic, while
the much larger C3 pool of PA may vary less in relative terms. Because the variance of the
contrast scales as the square of the coefficient of variation divided by the number of draws, a
10\% scatter needs about 100 draws per window, and the target is met by windows of feasible size
only if the scatter in PA is below about 10\%. A natural history study that sampled 50 patients
every six months for 36 months restricted its biomarker analysis to baseline, citing high
variability during follow-up~\cite{natural_history}.
% src: out/t42_serial_cv.txt:5-11 (0.133/46, 0.176/81, 0.475/588, 1.080/3033), :21-27, :29-38 (two-sided uncertainty of the anchor); REPORT.md:605,607 (no published within-patient CV in stable PA; Thompson 2012: 26.4%, 8.2-61.1%; natural history n = 50, 6-monthly to 36 months, analysis restricted to baseline), :649 (CV^2/n; ~100 at 10%); pi_note/pi_note.tex:297-303 (MMA serum MMA 31 +/- 25%, Manoli 2021; <~10%); paper/claims.csv:19 (C20, natural_history: visits every 6 months to month 36)

\paragraph{Drift in the patient constants.} The cancellation is exact only while $k_e$ and the
calibration of the load are the same in both windows. A drift in $\ln k_e$ of 0.03 between the
windows raises the contrast to 0.18, and a drift of 0.06 raises it to 0.28, about 42\um; the
contrast stays below 0.20 up to a drift of about 0.035. Over a week $k_e$, which includes the
renal clearance of the marker, is effectively fixed. Over a year in a growing child, whose renal
clearance scales with body size, or with declining kidney function, a drift of 0.05--0.10 is
plausible. A shift in the marker's calibration equal to 1.5\% of the load, as a change in
carnitine dose or in sampling conditions (fed or fasted state, time since the last dose) could
produce, raises the contrast to 0.204, about 30\um; the contrast stays below 0.20 up to a shift of
about 1.4\% of the load. Monitoring is therefore reliable over intervals of weeks to a few months
on a stable regimen with standardised sampling, and degrades over a year or across a change of
regimen.
% src: out/t30_monitoring_confounders.txt:9-29 (0.182 at 0.03, 0.280 at 0.06; 0.204 at 0.015 L); out/t24_contrast_stress.txt:15-36 (breakpoints ~0.035 and ~0.014 L; k_e fixed over a week; 0.05-0.1 over a year in a growing child); REPORT.md:596,1206 (0.28 -> 42 uM/d); pi_note/pi_note.tex:303-307 (42 uM/d; 30 uM/d)

\paragraph{Recorded changes in dose.} A change in the infused dose that is recorded enters the
model as a known change in $\Delta$ and does not mimic a change in $\theta$. The contrast stays
unbiased, at a cost in precision: 0.133 with the same dose in both windows and 0.184 if the dose
doubles in the second. An uncertain delivered dose would instead leak into $\hat\theta$, so the
date, dose and lot of every infusion must be recorded.
% src: out/t25_drug_cv_marker.txt:11-20 (same dose 0.133, double dose on window 2 0.184; uncertain dose leaks); REPORT.md:647

\paragraph{Unrecorded antibiotic courses.} Gut bacteria supply a substantial part of the load
(Supplementary Section~\ref{si:s1-bio}), and metronidazole lowered urinary propionate metabolites
by 41\% on average (12--76\%; Supplementary Section~\ref{si:s2:record}). A course that starts after
the first window and continues through the second lowers the load without the record showing it.
Read at the nominal recorded load, the lower marker values then look like higher capacity. In the
linear response, an unrecorded fall in load of 10\%, 20\% and 41\% shifts
$\ln\hat\theta_2-\ln\hat\theta_1$ by $+1.1$, $+2.0$ and $+2.9$, so that the course mimics a rise
of one to three times the patient's whole capacity, about 9 to 22 times the standard deviation of
the contrast (0.133). The shift is always towards apparent improvement.
Because antibiotics tend to be started when there is clinical concern, an unrecorded course would
make a deteriorating patient look better. Antibiotic status must therefore be recorded and held
fixed across the two windows. Within one window, an unrecorded antibiotic cycle also biases the
estimate of the level itself downwards (Supplementary Section~\ref{si:s2:record}).
% src: gut share (22.2 +/- 8.5%, params.py:496-505) and metronidazole fall (41%, 12-76%, claims.csv:50) cited in s1_model.tex (si:s1-bio) and s2_sampling.tex (si:s2:record); out/t24_contrast_stress.txt:38-50 (+1.14, +2.01, +2.94); REPORT.md:642 (always towards improvement; quoted as 8-20x the 0.13, recomputed 1.14/0.133 = 8.6, 2.94/0.133 = 22.1), :1206 (one to three times the whole capacity)

\begin{table}[htbp]
\centering\small\setlength{\tabcolsep}{5pt}
\caption{\label{stab:s5conditions}Conditions on the within-patient contrast at 5\% of normal,
linear-clearance model, two windows of 104 draws. Contrast: Fisher standard deviation of
$\ln\theta_2-\ln\theta_1$, equal to $\operatorname{sd}(\hat\theta_2-\hat\theta_1)/150$; the
baseline, with a 5\% per-draw scatter and the constants exactly shared, is 0.133 (20\um). Draws:
draws per window needed for a contrast of 0.20. Per-draw scatter: coefficient of variation of each
marker value about the model's prediction; 8.2, 26.4 and 61.1\% are the range and mean of the
within-day variability of serum C3 in healthy adults~\cite{thompson2012c3}. Drifts are changes
between the windows that the fit treats as shared. The antibiotic row gives a bias rather than a
standard deviation: the linear-response shift of $\ln\hat\theta_2-\ln\hat\theta_1$, which is also the
apparent rise as a multiple of $\theta$.}
\begin{tabular}{@{}llcc@{}}
\toprule
condition & value & contrast & draws \\
\midrule
per-draw scatter of the marker & 5\% (assay) & 0.133 & 46 \\
& 8.2\% & 0.176 & 81 \\
& 15\% & -- & 208 \\
& 26.4\% & 0.475 & 588 \\
& 40\% & -- & 1315 \\
& 61.1\% & 1.080 & 3033 \\
\addlinespace
drift in $\ln k_e$ between windows & 0.02 & 0.157 & \\
& 0.03 & 0.182 & \\
& 0.05 & 0.245 & \\
& 0.06 & 0.280 & \\
& 0.10 & 0.426 & \\
& 0.20 & 0.767 & \\
\addlinespace
shift in calibration (fraction of load) & 0.01 & 0.176 & \\
& 0.015 & 0.204 & \\
& 0.02 & 0.226 & \\
& 0.03 & 0.254 & \\
& 0.05 & 0.276 & \\
\addlinespace
recorded dose doubled in window 2 & & 0.184 & \\
\addlinespace
unrecorded antibiotic course through window 2 & load fall 10\% & bias $+1.14$ & \\
(linear response) & load fall 20\% & bias $+2.01$ & \\
& load fall 41\% & bias $+2.94$ & \\
\bottomrule
\end{tabular}
\end{table}
% src: out/t42_serial_cv.txt:7-11,21-27; out/t24_contrast_stress.txt:16-28,38-45; out/t30_monitoring_confounders.txt:9-25; out/t25_drug_cv_marker.txt:12

% -----------------------------------------------------------------------------
\subsection{Other instruments and the saturable-clearance model}
\label{si:s5:other}

Measuring the shared constants does not improve the contrast, because they already cancel: the
intravenous tracer and the load turnover of Supplementary Section~\ref{si:s4:other} tighten what
the difference removes. An ex vivo PCC assay at 20\% repeatability in each window would give a
contrast of $\sqrt2\times0.20\approx0.28$, worse than the record's 0.13. A direct reading in vivo,
such as the breath test, is different: its obstacle for the level, a recovery that varies between
patients, cancels in the ratio of two readings from one patient. In our model its signal, the
fraction of the load oxidised, changes by 0.89\% for a 1\% change in capacity at the untreated
steady state at 5\% of normal, so at an assumed within-subject coefficient of variation of 8\% one
reading per window would give a contrast of 0.127, comparable to 104 blood draws. The calculation
does not include the delivered enzyme, which on therapy carries most of the load
(Supplementary Section~\ref{si:s4:split}). No within-subject coefficient of variation of
the breath test has been published for PA, and the recovery would itself have to stay constant
between the windows, although it may change with diet, gut flora or antibiotics. The value is
therefore a design target, which a measurement of the test's within-subject repeatability in PA
would have to confirm.
% src: out/t26_multiobservable.txt:42-47 (sensors add nothing; ex vivo sqrt2 x 0.20 ~ 0.28); out/t31_breath_contrast.txt:5-18 (untreated operating point L 322.7, m 54.0, S 0.163; elasticity 0.888; 8% one reading 0.127), :35-43 (CV assumed, not measured); pi_note/figure_guide.tex:272-276 (own enzyme about 6% of the load on therapy); REPORT.md:618-631 (CV_S a sweep; c drifts with diet, gut flora, antibiotic course), :731

In the saturable-clearance model, where the enzyme carries a larger share of the load, the
Fisher value of the contrast is 0.088 at the nominal constants and 0.078 at the constructed
operating point of Supplementary Section~\ref{si:s4:sat}, against 0.133 in the linear-clearance
model on the same footing (0.132 at the managed load of 220\um). The shared constants, now
including the maximal rate $V_e$ and half-saturation $K_e$ of the non-enzymatic route, again
shift both windows' estimates by an amount. We did not simulate pairs in that model, so the
reading of the change as a difference rests there on the Fisher calculation alone.
% src: out/t33_model_b.txt:56-64 (0.133 vs 0.088; shared Km, b, ke, Ke); out/t35_managed_load.txt:41-50 (0.132 vs 0.078); out/t36_reconcile_remark15.txt:54-62 (b shared; B better on both axes); manuscript/main.tex:532-533 (simulated in the linear-clearance model only)

% -----------------------------------------------------------------------------
\subsection{What would test the result}
\label{si:s5:test}

The contrast rests on the shared constants cancelling, which is a property of the model. Its
checkable consequence is the marker's own serial variability within a stable patient. With draws
taken at a fixed point of the dose cycle, the model predicts a within-patient standard deviation
of $\ln m$ of about 0.12 (a coefficient of variation of about 12\%), from the 10\% day-to-day
variation in intake (0.099 alone) and the 5\% assay (0.050 alone); with draws at random points of
the cycle it predicts about 0.64, dominated by the cycle. If stable patients sampled at a fixed
point of the cycle vary materially more than 12\%, either $\theta$ is not constant between visits
or a nuisance does not cancel, and in either case the contrast is optimistic. Measuring this
serial variability in PA, with the sampling time relative to the infusion, fed state, carnitine
dose and renal function recorded at each draw, is the measurement that monitoring needs first
(main text, Table~2).
% src: out/t25_drug_cv_marker.txt:44-58 (0.636 random phase; 0.121 synchronised; load 0.099, assay 0.050); REPORT.md:651; out/t30_monitoring_confounders.txt:41-52; manuscript/main.tex:383,564-565 (measured first; Table 2, monitoring row)
